\documentclass[a4paper,11pt]{article}

\usepackage[utf8]{inputenc}
\usepackage[T1]{fontenc}
\usepackage[margin=2.5cm]{geometry}
\usepackage{amsmath}
\usepackage{listings}
\usepackage{xcolor}

\lstset{
    language=C++,
    basicstyle=\ttfamily\small,
    keywordstyle=\color{blue!60!black},
    commentstyle=\color{green!40!black},
    numbers=left,
    numberstyle=\tiny,
    frame=single,
    breaklines=true,
    showstringspaces=false
}

\title{WMCS045-05 -- Lab Assignment \#4}
\author{Jort Oosterveld}
\date{2026-09-30}

\begin{document}

\maketitle

\section{Squared Sums}

\subsection{System}

All measurements were taken on the system in Table~\ref{tab:system}. The
benchmark was pinned to one core with \texttt{taskset}.

\begin{table}[ht]
    \centering
    \caption{System used for the measurements.}
    \label{tab:system}
    \begin{tabular}{ll}
        \hline
        CPU & AMD Ryzen 7 5825U (Zen 3, family 25, model 80) \\
        Cores / threads & 8 / 16 \\
        Frequency & max. 4.55\,GHz (boost enabled), 4.50\,GHz assumed for cycle counts \\
        Caches & L1d 32\,KiB and L2 512\,KiB per core, L3 16\,MiB shared \\
        SIMD registers & 16 (AVX2, FMA3) \\
        OS & Linux 7.0.0-31-generic \\
        Governor & \texttt{performance} (\texttt{amd-pstate-epp}) \\
        Compiler & g++ (GCC) 15.2.1 \\
        Flags & \texttt{-std=c++17 -O2 -march=native} \\
              & \texttt{-fno-tree-vectorize -ffp-contract=off} \\
        \hline
    \end{tabular}
\end{table}

The microarchitecture parameters of Zen 3 that matter for this exercise are
listed in Table~\ref{tab:uarch}. They are taken from the AMD Software
Optimization Guide for family 19h processors and the instruction tables of
uops.info.

\begin{table}[ht]
    \centering
    \caption{Relevant microarchitecture parameters of Zen 3 (scalar double precision).}
    \label{tab:uarch}
    \begin{tabular}{lccc}
        \hline
        Instruction & Latency (cycles) & Throughput (per cycle) & Execution units \\
        \hline
        \texttt{vaddsd} & 3 & 2 & FP2, FP3 \\
        \texttt{vmulsd} & 3 & 2 & FP0, FP1 \\
        \texttt{vfmadd231sd} & 4 & 2 & FP0, FP1 \\
        \hline
        \multicolumn{4}{l}{Dispatch width: 6 macro-ops per cycle} \\
        \hline
    \end{tabular}
\end{table}

\subsection{Benchmark}

The loop of \texttt{sqsum} is unrolled by a factor $U$ with one accumulator per unrolled statement, which breaks the dependency for the one accumulator.
The complete benchmark program \texttt{loop\_unrolling.cpp} can be found in Listing~\ref{lst:full} in Appendix~\ref{app:code}.

Two compiler flags are needed to measure what is intended.
\texttt{-fno-tree-vectorize} keeps the loop scalar, and \texttt{-ffp-contract=off} keeps the separate multiply and add, otherwise GCC fuses them into an FMA by itself.

As in the assignment, $n = 8192$ and \texttt{sqsum} is called 10000 times.
This batch of 10000 calls is timed with
\texttt{clock\_gettime(CLOCK\_MONOTONIC\_RAW)} and repeated 30 times per
unrolling factor. Only the minimum (best case) of the 30 repetitions is reported, since
disturbances by the operating system can only make a run slower.

Hardware cycle counters are not accessible on this system
(\texttt{perf\_event\_paranoid} = 4). The time per element is therefore
converted to cycles with an assumed constant core frequency of 4.50\,GHz,
the boost clock the core runs at with the \texttt{performance} governor
while the benchmark is pinned to it. The cycle counts are therefore only as
accurate as this assumption, the times are what was measured.

\subsection{Results}

Table~\ref{tab:results} shows the results (\texttt{loop\_unrolling.csv}).

\begin{table}[ht]
    \centering
    \caption{Best case of 30 repetitions for the squared sum with unrolling
    factor $U$, $n = 8192$, 10000 calls per repetition. The speed-up is
    relative to the loop without unrolling.}
    \label{tab:results}
    \begin{tabular}{rcccc}
        \hline
        $U$ & $\mu$s / call & ns / element & cycles / element & speed-up \\
        \hline
        1  & 5.49 & 0.6705 & 3.017 & 1.00 \\
        2  & 2.78 & 0.3388 & 1.525 & 1.98 \\
        3  & 1.88 & 0.2297 & 1.034 & 2.92 \\
        4  & 1.55 & 0.1888 & 0.850 & 3.55 \\
        5  & 1.25 & 0.1532 & 0.689 & 4.38 \\
        6  & 1.23 & 0.1498 & 0.674 & 4.48 \\
        7  & 1.20 & 0.1464 & 0.659 & 4.58 \\
        8  & 1.07 & 0.1306 & 0.587 & 5.13 \\
        9  & 1.14 & 0.1390 & 0.625 & 4.82 \\
        10 & 1.12 & 0.1365 & 0.614 & 4.91 \\
        11 & 1.06 & 0.1300 & 0.585 & 5.16 \\
        12 & 1.15 & 0.1404 & 0.632 & 4.78 \\
        14 & 1.05 & 0.1280 & 0.576 & 5.24 \\
        16 & 1.62 & 0.1982 & 0.892 & 3.38 \\
        \hline
    \end{tabular}
\end{table}

\subsection{a) Optimal unrolling factor}

Without unrolling, every addition has to wait for the previous one, so one
element costs the latency of the loop-carried operation. Only the addition is
part of this dependency chain, the multiplication depends on the loaded value
alone. The measured 3.02 cycles per element for $U = 1$ match the latency of
3 cycles of \texttt{vaddsd}.

With $U$ accumulators there are $U$ independent chains, so the cost drops to
$3 / U$ cycles per element until the throughput of the core becomes the
limit. Zen 3 executes 2 additions and 2 multiplications per cycle, which is
0.5 cycles per element. The chains stop being the bottleneck at
\[
    U_\text{opt} = \text{latency} \times \text{throughput} = 3 \times 2 = 6 .
\]

The measurements follow $3 / U$ up to $U = 3$ and then flatten out.
after $U = 6$, there is no systematic gain anymore.
All factors from 6 to 14 lie between 0.58 and 0.67 cycles per element.
This agrees with the expected factor of 6, which is 4.5 times faster than the original loop.
The theoretical 0.5 cycles per element is not reached.
A likely reason is the dispatch width: the loop needs 3 macro-ops per element, a load, multiply and add.
So 0.5 cycles per element would already use all 6 macro-ops per cycle that Zen 3 can dispatch, and the loop overhead comes on top.
This also explains why $U = 8$ is clearly faster than $U = 6$ with 0.587 cycles per element.
There is less loop overhead (shown in listings~\ref{lst:u6_asm} and~\ref{lst:u8_asm} when unrolling is done with a power of two, which gives an estimate of $(3 \cdot 8 + 2) / (6 \cdot 8) = 0.54$ cycles per element against $(3 \cdot 6 + 4) / (6 \cdot 6) = 0.61$ for $U = 6$.
With the same estimate the overhead of 4 macro-ops is spread over more elements for
larger factors, so $U = 11$ and $U = 14$ are as fast as $U = 8$.
The differences between the factors 8 to 14 can be attributed to variations between runs.

\begin{lstlisting}[caption={Loop overhead of compiled assembler with u=6}, label=lst:u6_asm]
    leal   11(%rcx), %edx     
    addq   $48, %rax          
    addl   $6, %ecx           
    cmpl   %edx, %esi         
    jg     .L51                
\end{lstlisting}

\begin{lstlisting}[caption={Loop overhead of compiled assembler with u=8}, label=lst:u8_asm]
    addq   $64, %rax          
    cmpq   %rax, %rdx         
    jne    .L80                 
\end{lstlisting}

\subsection{b) Fused-multiply-add}

With FMA, the multiplication and addition are one instruction, so there are
fewer 2 macro-ops per element instead of 3. The drawback is that the whole 4 cycles of FMA are dependent on each other, instead of the 3 from the addition.
Without unrolling, FMA is therefore expected to be about $\frac{1}{3}$ slower.

The throughput of FMA on Zen 3 is 2 per cycle, the same as for the separate
multiply and add.
But, more accumulators are needed to reach it, $4 \times 2 = 8$ instead of 6.
For every unrolling factor below 8 the dependency chain of the FMA is still the bottleneck, so FMA is slower.

FMA can only be better when the loop is unrolled by at least 8.
The execution units then allow the same 0.5 cycles per element for both
variants, but the multiply and add loop needs all 6 macro ops per cycle of
the dispatch width to get there, which is the likely reason that it only
reached 0.58 cycles per element in the measurements. With 2 macro ops per
element FMA needs only 4 per cycle, so it could come closer to the limit.
\\\\The relevant microarchitecture parameters are:
\begin{itemize}
    \item the latency of the addition compared to the latency of the FMA,
    since this is the length of the loop-carried dependency chain,
    \item the throughput (number of execution units) of addition,
    multiplication and FMA, which sets the limit once the dependency is
    broken,
    \item the pipeline (dispatch) width in relation to the number of
    macro-ops per element,
    \item the number of registers, which limits the number of accumulators
    (15 on this system).
\end{itemize}

\appendix

\section{Benchmark source code}
\label{app:code}

Listing~\ref{lst:full} is the complete benchmark program
\texttt{loop\_unrolling.cpp} used for question 1. It is compiled and run
with \texttt{run\_benchmark.sh}, which pins it to one core with
\texttt{taskset} and writes \texttt{loop\_unrolling.csv}.

\lstinputlisting[caption={Complete benchmark program \texttt{loop\_unrolling.cpp}.}, label={lst:full}]{../loop_unrolling.cpp}

\end{document}
